\subsection{What the Reading Operator Itself Contributes}
\label{sec:operator}

The floor reported above is the error of a pair of operators: the
Clough--Tocher transfer that places the data on the carrier, and the
nearest-cell read that retrieves them. Appendix~\ref{app:operator}
characterises the read alone, on analytic single-frequency fields placed on
the cell centres of three regions with nothing resampled. The read is
registered to well under a cell (largest fitted displacement
\SI{2.1}{\milli\metre}), its amplitude response is $0.99$ or higher beyond
\SI{100}{m}, and almost all of its error ($95\%$ at \SI{21}{m}, essentially
$100\%$ above \SI{60}{m}) is a residual that no amplitude-and-phase model
describes. Reading the containing cell is therefore not a filter but a
sub-pixel positional jitter of about \SI{2.1}{m} on a \SI{10}{m} native
pixel, and this jitter is what the floor is largely made of. On a
band-limited field carrying the spectrum of the imagery ($0.19\%$ of its
variance above half the native Nyquist frequency, against the $0.21\%$
median measured on the crops in Section~\ref{sec:support_controls}), the
read alone produces an RMSE of $14.0$--$14.4\%$ of the standard deviation of
the field, against a floor of about $10\%$ on the real tiles. The two
figures are consistent: a carrier tile is a Clough--Tocher interpolation of
the scene rather than a point sampling of it, and the smoothing removes the
gradients that a jitter converts into error.

The floor of Table~\ref{tab:representation_floor} therefore cannot be read
as the representation error of the carrier grid itself. It is a property of
the transfer-and-read pair and an upper bound on the cost of the
representation. What would lower it is not a better low-pass filter, since
the amplitude response is already above $0.99$ over the band that carries
the variance, but a reconstructor that interpolates within the carrier
lattice instead of reading the containing cell. Such a reconstructor was
implemented, and Table~\ref{tab:reader_floor} reports what it removes.

\begin{table}[t]
\centering
\caption{The transfer-and-read floor under three readers}
\label{tab:reader_floor}
\footnotesize
\setlength{\tabcolsep}{4pt}
\begin{tabular}{@{}llrrrr@{}}
\toprule
Tiles & Reader & $n$ & Floor & \% of std & Cells read \\
\midrule
Clough--Tocher & nearest & 40 & 4.04e-03 & 10.0 & 30.8\% \\
Clough--Tocher & bilinear & 40 & 1.94e-03 & 5.9 & 78.4\% \\
PSF & nearest & 40 & 7.83e-03 & 23.6 & 30.8\% \\
PSF & bilinear & 40 & 4.29e-03 & 13.8 & 78.4\% \\
PSF & psf & 40 & 4.44e-03 & 13.2 & 81.1\% \\
\midrule
\multicolumn{2}{@{}l}{Clough--Tocher: bilinear / nearest} & 40 & 0.582 & \multicolumn{2}{l@{}}{MSE reduction, 66\%} \\
\multicolumn{2}{@{}l}{PSF: bilinear / nearest} & 40 & 0.581 & \multicolumn{2}{l@{}}{MSE reduction, 66\%} \\
\bottomrule
\end{tabular}
\end{table}


The interpolating read places each native pixel in the lattice of the
carrier at sub-cell precision. The local index of a NESTED block is the
Morton code of the integer position of the cell on its base face, so the
tile is a $1024\times1024$ lattice whose neighbours are spatial neighbours.
The read inverts, at the containing cell, the affine map from lattice steps
to metres measured on the cell centres, and takes the bilinear combination
of the four bracketing centres; anisotropy and orientation enter through
this Jacobian, the weights sum to one, and the read is linear. On the same
tiles, mask and observation, its floor is $0.582$ times that of the nearest
read (paired median over the 40 regions, lower on every region), so that
the read accounts for $66\%$ of the variance of the floor: the
transfer-and-read floor decreases from $10.0\%$ to $5.9\%$ of the standard
deviation of a scene when the read is changed and nothing else. On the
PSF-transferred tiles the same read lies within $4\%$ of the forward
operator of the companion study ($0.957$). The floor is thus two thirds a
property of how the carrier is read and one third a property of what was
placed on it, and every result of Section~\ref{sec:common_reference} scored
through the nearest read inherits this decomposition.

The read has a second consequence, which concerns rate rather than error. A
nearest-cell read of a $1024^2$ tile by the $637^2$ crop touches only
$30.8\%$ of the cells (median over 40 regions, on the intersection mask);
the remaining cells form the kernel of the operator, and whatever a codec
writes there never reaches the product. Using as error source a uniform
dead-zone quantiser in an orthonormal wavelet domain, whose error is close
to white in cell space, and bisecting its step onto the RMSE achieved by
each codec, $68.9\%$ of the error energy lies in cells that the product
never reads. This does not make the product error smaller: the RMS of the
transmitted error is $1.00$ times its RMS over all cells, because a gather
leaves the error per read pixel unchanged, so a coarser step gains nothing
($0.1\%$ of the rate). What the kernel reveals is allocation, namely bits
spent on cells that nothing depends on, with a bound of
$1/0.308\approx3.2\times$ that a uniform quantiser cannot reach. Seen from
the sweep, $72.5\%$ of the $2\,080$ codec operating points reach a
cell-space RMSE below the floor of Table~\ref{tab:representation_floor}.

The matched case behaves differently. On the PSF-transferred tiles of
Section~\ref{sec:transfer_cost}, read through the forward operator for
which they were built (floor $0.558$ times that of the nearest read, lower
on all 40 regions), the operator attenuates a white coding error to $0.67$
of its cell-space RMS, and setting the distortion target on the product
rather than on the carrier reaches the same product distortion with $20\%$
fewer bytes (median; $4.6\%$ at the tightest quartile of targets, $36\%$ at
the loosest). Three qualifications prevent this from being a
recommendation. It is a property of the reader rather than of the grid,
since a gather has a large kernel precisely because the carrier is
oversampled relative to it; white error is the most favourable case for
attenuation; and the gain holds only for a product consumed through that
operator, since noise shaped into its stop-band remains visible to anyone
who reads the carrier as an image. What the measurement establishes is
narrower: the distortion that a carrier-grid study should target is the one
its reader transmits, and on this carrier the two are not the same
quantity.
